% Generated draft: requires mathematical and source review.
\documentclass[11pt]{article}
\usepackage[margin=25mm]{geometry}
\usepackage{fontspec}
\setmainfont{Latin Modern Roman}
\usepackage{amsmath,amssymb,mathtools,unicode-math}
\setmathfont{Latin Modern Math}
\usepackage{enumitem,needspace,xurl,hyperref}
\setlength{\parindent}{0pt}
\setlength{\parskip}{5pt}
\setlength{\emergencystretch}{3em}
\newcommand{\N}{\mathbb{N}}
\newcommand{\Z}{\mathbb{Z}}
\newcommand{\Q}{\mathbb{Q}}
\newcommand{\R}{\mathbb{R}}
\newcommand{\C}{\mathbb{C}}
\newcommand{\F}{\mathbb{F}}
\newcommand{\A}{\mathbb{A}}
\newcommand{\PP}{\mathbb{P}}
\newcommand{\E}{\mathbb{E}}
\newcommand{\eps}{\varepsilon}
\newcommand{\dd}{\,\mathrm{d}}
\begin{document}
\section*{UnsolvedMath ID: 30003357}
\textbf{Source problem code:} catalogue entry
\subsection{Definitions and mathematical statement}
Let \(M\) be a connected smooth manifold, and let \(\pi\colon\widetilde M\to M\) denote its universal covering manifold. A manifold is called non-spin here if its tangent bundle admits no spin structure; the supplied record does not clarify whether orientability is an additional intended hypothesis. Write
\[
\mathcal R(M)=\{g:g\text{ is a smooth Riemannian metric on }M\},
\]
and define
\[
\mathcal R_{\operatorname{Ric}>0}(M)
=\{g\in\mathcal R(M):\operatorname{Ric}_g(v,v)>0
\text{ for every }x\in M\text{ and }0\ne v\in T_xM\}.
\]
For the positive-scalar-curvature spaces mentioned in the supplied background, write
\[
\mathcal R_{\operatorname{scal}>0}(M)
=\{g\in\mathcal R(M):\operatorname{scal}_g(x)>0
\text{ for every }x\in M\}.
\]
\textbf{Source question.} For manifolds \(M\) satisfying either
\[
M\text{ is simply connected and non-spin},
\qquad\text{or}\qquad
\widetilde M\text{ is non-spin},
\]
what can be determined about their spaces of Riemannian metrics, especially \(\mathcal R_{\operatorname{Ric}>0}(M)\)?

This is a research question rather than a specified conjecture. The record does not identify a particular invariant to determine, assert that positive-Ricci metrics exist, or specify a topology on these sets. It also does not specify dimension, compactness, completeness, boundary conditions, or whether metrics are to be considered modulo diffeomorphisms. The displayed definitions describe sets of metrics and impose none of these additional requirements. A precise topological or moduli-space problem therefore requires further hypotheses and a specified target.

The supplied dated curation labels the problem partially solved. It attributes to Ebert and Wiemeler an identification of the homotopy type of the positive-scalar-curvature space for simply connected non-spin manifolds in dimensions at least five other than eight, while describing the broader non-spin and positive-Ricci questions as open. This is a supplied literature assessment, not an independently verified theorem statement; any additional hypotheses of the cited result must be checked before using it.
\subsection{Short English statement}
The question asks what can be learned about spaces of Riemannian metrics, especially metrics with positive Ricci curvature, when the manifold is simply connected and non-spin or its universal cover is non-spin. The source leaves the precise invariant, topology, and several manifold hypotheses unspecified. Its partially solved label and positive-scalar-curvature progress report are supplied dated curation requiring verification.
\subsection{Sources}
\begin{itemize}
\item Local UnsolvedMath registry: supplied complete record with permanent ID \(30003357\), source problem code \texttt{catalogue entry}, title ``Positive Ricci Metric Spaces on Nonspin Manifolds,'' in the Oberwolfach Reports Open Problems collection. No individual proposer is named. The registry marks the statement as exact and describes it as retained byte-for-byte; this is supplied internal verification information.
\item Supplied workshop citation: \emph{Mini-Workshop: Spaces and Moduli Spaces of Riemannian Metrics (2018)}, Oberwolfach Reports, pp.\ 133--166, DOI \texttt{10.4171/owr/2017/3}. \url{https://doi.org/10.4171/owr/2017/3}. The record attributes the question to collected workshop open-problem discussions. Its proposed year and year tag are 2017, whereas the supplied report title says 2018; this discrepancy remains unresolved.
\item Supplied literature source, ``Ebert--Wiemeler on positive scalar curvature spaces'': \url{https://doi.org/10.4171/JEMS/1333}.
\item Supplied literature source, ``Reiser on positive Ricci curvature'': \url{https://doi.org/10.1112/topo.70007}.
\item Supplied literature source, ``Oberwolfach report on non-spin metrics'': \url{https://doi.org/10.4171/owr/2017/3}.
\item Literature assessment checked at \texttt{2026-08-21}: the supplied curation labels the problem partially solved and reports the positive-scalar-curvature progress described above, with broader non-spin and positive-Ricci questions remaining open. This date and assessment are supplied dated curation derived from the supplied OWR curation, not independently verified source visits or a new literature assessment. The supplied background explicitly requests independent verification of time-sensitive status and cited claims.
\end{itemize}
\end{document}
